% answers_short.tex —— 1982 年·全国卷（文科） · 简短答案速查
% 由 tools/make_short_answers.py 自 answers.tex 的 \gksolhead/\gkanswer 提取，
% 并与 1982/1982_数学真题_文_答案要点.md 逐题比对后生成，勿手改（改源 MD 后重跑）。

\begin{gknotice}
\textbf{使用说明}：本节只给最终答案便于快速核对，完整推导见解析版。答案为编写者独立推导，\textbf{非官方标准答案}。
\end{gknotice}

\gkpart{题目}
\begin{itemize}[leftmargin=2.6em]
  \item[\textbf{1.}] 逐行依次为 $x=0$；全体实数 $\mathbf{R}$；$x>0$；全体实数 $\mathbf{R}$
  \item[\textbf{2.}] $-38760$
  \item[\textbf{3.}] $4x^2+y^2=4$ 表示椭圆 $\dfrac{x^2}{1}+\dfrac{y^2}{4}=1$（中心在原点，长轴在 $y$ 轴上，顶点为 $(\pm1,0)$、$(0,\pm2)$）；$x-3=0$ 表示垂直于 $x$ 轴且过点 $(3,0)$ 的直线
  \item[\textbf{4.}] $\pm\dfrac{\sqrt7}{4}$
  \item[\textbf{5.}] 当长（平行于墙的一边）为 $\dfrac L2$、宽（垂直于墙的一边）为 $\dfrac L6$ 时场地面积最大，最大面积为 $\dfrac{L^2}{12}$
  \item[\textbf{6.}] (1) $\dfrac{5a^3}{6}$；(2) $60^\circ$
  \item[\textbf{7.}] 取 $AB$ 所在直线为 $x$ 轴、$AB$ 中点为原点，则 $A(-a,0)$、$B(a,0)$，轨迹方程为 $\left(x-\dfrac{5a}{3}\right)^2+y^2=\dfrac{16a^2}{9}$（即 $3x^2+3y^2-10ax+3a^2=0$），它表示以 $\left(\dfrac{5a}{3},0\right)$ 为圆心、$\dfrac{4a}{3}$ 为半径的圆（阿波罗尼斯圆）
  \item[\textbf{8.}] $4$
  \item[\textbf{9.}] $\dfrac{a^2b\sin\theta}{2(a+b\cos\theta)}$
\end{itemize}
